% Options for packages loaded elsewhere
\PassOptionsToPackage{unicode}{hyperref}
\PassOptionsToPackage{hyphens}{url}
%
\documentclass[
]{article}
\usepackage{amsmath,amssymb}
\usepackage{iftex}
\ifPDFTeX
  \usepackage[T1]{fontenc}
  \usepackage[utf8]{inputenc}
  \usepackage{textcomp} % provide euro and other symbols
\else % if luatex or xetex
  \usepackage{unicode-math} % this also loads fontspec
  \defaultfontfeatures{Scale=MatchLowercase}
  \defaultfontfeatures[\rmfamily]{Ligatures=TeX,Scale=1}
\fi
\usepackage{lmodern}
\ifPDFTeX\else
  % xetex/luatex font selection
\fi
% Use upquote if available, for straight quotes in verbatim environments
\IfFileExists{upquote.sty}{\usepackage{upquote}}{}
\IfFileExists{microtype.sty}{% use microtype if available
  \usepackage[]{microtype}
  \UseMicrotypeSet[protrusion]{basicmath} % disable protrusion for tt fonts
}{}
\makeatletter
\@ifundefined{KOMAClassName}{% if non-KOMA class
  \IfFileExists{parskip.sty}{%
    \usepackage{parskip}
  }{% else
    \setlength{\parindent}{0pt}
    \setlength{\parskip}{6pt plus 2pt minus 1pt}}
}{% if KOMA class
  \KOMAoptions{parskip=half}}
\makeatother
\usepackage{xcolor}
\usepackage{color}
\usepackage{fancyvrb}
\newcommand{\VerbBar}{|}
\newcommand{\VERB}{\Verb[commandchars=\\\{\}]}
\DefineVerbatimEnvironment{Highlighting}{Verbatim}{commandchars=\\\{\}}
% Add ',fontsize=\small' for more characters per line
\newenvironment{Shaded}{}{}
\newcommand{\AlertTok}[1]{\textcolor[rgb]{1.00,0.00,0.00}{\textbf{#1}}}
\newcommand{\AnnotationTok}[1]{\textcolor[rgb]{0.38,0.63,0.69}{\textbf{\textit{#1}}}}
\newcommand{\AttributeTok}[1]{\textcolor[rgb]{0.49,0.56,0.16}{#1}}
\newcommand{\BaseNTok}[1]{\textcolor[rgb]{0.25,0.63,0.44}{#1}}
\newcommand{\BuiltInTok}[1]{\textcolor[rgb]{0.00,0.50,0.00}{#1}}
\newcommand{\CharTok}[1]{\textcolor[rgb]{0.25,0.44,0.63}{#1}}
\newcommand{\CommentTok}[1]{\textcolor[rgb]{0.38,0.63,0.69}{\textit{#1}}}
\newcommand{\CommentVarTok}[1]{\textcolor[rgb]{0.38,0.63,0.69}{\textbf{\textit{#1}}}}
\newcommand{\ConstantTok}[1]{\textcolor[rgb]{0.53,0.00,0.00}{#1}}
\newcommand{\ControlFlowTok}[1]{\textcolor[rgb]{0.00,0.44,0.13}{\textbf{#1}}}
\newcommand{\DataTypeTok}[1]{\textcolor[rgb]{0.56,0.13,0.00}{#1}}
\newcommand{\DecValTok}[1]{\textcolor[rgb]{0.25,0.63,0.44}{#1}}
\newcommand{\DocumentationTok}[1]{\textcolor[rgb]{0.73,0.13,0.13}{\textit{#1}}}
\newcommand{\ErrorTok}[1]{\textcolor[rgb]{1.00,0.00,0.00}{\textbf{#1}}}
\newcommand{\ExtensionTok}[1]{#1}
\newcommand{\FloatTok}[1]{\textcolor[rgb]{0.25,0.63,0.44}{#1}}
\newcommand{\FunctionTok}[1]{\textcolor[rgb]{0.02,0.16,0.49}{#1}}
\newcommand{\ImportTok}[1]{\textcolor[rgb]{0.00,0.50,0.00}{\textbf{#1}}}
\newcommand{\InformationTok}[1]{\textcolor[rgb]{0.38,0.63,0.69}{\textbf{\textit{#1}}}}
\newcommand{\KeywordTok}[1]{\textcolor[rgb]{0.00,0.44,0.13}{\textbf{#1}}}
\newcommand{\NormalTok}[1]{#1}
\newcommand{\OperatorTok}[1]{\textcolor[rgb]{0.40,0.40,0.40}{#1}}
\newcommand{\OtherTok}[1]{\textcolor[rgb]{0.00,0.44,0.13}{#1}}
\newcommand{\PreprocessorTok}[1]{\textcolor[rgb]{0.74,0.48,0.00}{#1}}
\newcommand{\RegionMarkerTok}[1]{#1}
\newcommand{\SpecialCharTok}[1]{\textcolor[rgb]{0.25,0.44,0.63}{#1}}
\newcommand{\SpecialStringTok}[1]{\textcolor[rgb]{0.73,0.40,0.53}{#1}}
\newcommand{\StringTok}[1]{\textcolor[rgb]{0.25,0.44,0.63}{#1}}
\newcommand{\VariableTok}[1]{\textcolor[rgb]{0.10,0.09,0.49}{#1}}
\newcommand{\VerbatimStringTok}[1]{\textcolor[rgb]{0.25,0.44,0.63}{#1}}
\newcommand{\WarningTok}[1]{\textcolor[rgb]{0.38,0.63,0.69}{\textbf{\textit{#1}}}}
\setlength{\emergencystretch}{3em} % prevent overfull lines
\providecommand{\tightlist}{%
  \setlength{\itemsep}{0pt}\setlength{\parskip}{0pt}}
\setcounter{secnumdepth}{-\maxdimen} % remove section numbering
\ifLuaTeX
  \usepackage{selnolig}  % disable illegal ligatures
\fi
\IfFileExists{bookmark.sty}{\usepackage{bookmark}}{\usepackage{hyperref}}
\IfFileExists{xurl.sty}{\usepackage{xurl}}{} % add URL line breaks if available
\urlstyle{same}
\hypersetup{
  hidelinks,
  pdfcreator={LaTeX via pandoc}}

\author{}
\date{}

\begin{document}

\hypertarget{swaps}{%
\section{Swaps}\label{swaps}}

This guide walks you through testing Ark/Lightning \textbf{submarine}
and \textbf{reverse submarine} swaps against a full local regtest stack.

The stack is orchestrated end-to-end by the in-house
\texttt{arkade-regtest} Node CLI (\texttt{regtest/regtest.mjs}, vendored
as the \texttt{regtest} git submodule) --- there is no Nigiri and no
manual service wiring. A single \texttt{pnpm\ run\ regtest:start} brings
up \textbf{and provisions} everything:

\begin{itemize}
\tightlist
\item
  \textbf{Ark} --- Bitcoin Core + \texttt{arkd}/\texttt{arkd-wallet}
  (the server wallet is created, unlocked and funded automatically).
\item
  \textbf{Boltz} --- the Boltz backend, its Lightning node
  (\texttt{boltz-lnd}) and Fulmine (\texttt{boltz-fulmine}). The
  \texttt{boltz-lnd\ ⇄\ lnd} channel is opened and balanced, Boltz's
  Bitcoin Core wallet is funded, and the ARK/BTC pairs are verified.
\item
  \textbf{Counterparty Lightning} --- a second LND node (\texttt{lnd})
  that you drive with \texttt{lncli} to play the remote side of each
  swap.
\item
  \textbf{Arkade wallet} --- the app under test (run with
  \texttt{pnpm\ start}).
\end{itemize}

\begin{quote}
Everything the earlier (Nigiri-based) edition of this guide did by hand
--- \texttt{nigiri\ start\ -\/-ln}, manual \texttt{arkd\ wallet}
create/unlock, opening LND channels, wiring Fulmine --- is now
automatic. Funding is done with the CLI faucet:
\texttt{node\ regtest/regtest.mjs\ faucet\ \textless{}address\textgreater{}\ \textless{}btc\textgreater{}\ -\/-confirm}.
\end{quote}

\hypertarget{requirements}{%
\subsection{Requirements}\label{requirements}}

\begin{itemize}
\tightlist
\item
  \href{https://docs.docker.com/engine/install/}{Docker}
\item
  \href{https://nodejs.org/}{Node.js} v20.19+ or v22.12+ (drives the
  \texttt{arkade-regtest} stack via \texttt{regtest/regtest.mjs})
\item
  \href{https://formulae.brew.sh/formula/jq}{jq} (optional --- used
  below to pull fields out of \texttt{lncli} JSON)
\end{itemize}

\hypertarget{start-the-regtest-stack}{%
\subsection{1. Start the regtest stack}\label{start-the-regtest-stack}}

From the repo root:

\begin{Shaded}
\begin{Highlighting}[]
\ExtensionTok{pnpm}\NormalTok{ run regtest:start}
\end{Highlighting}
\end{Shaded}

This runs
\texttt{node\ regtest/regtest.mjs\ start\ -\/-env\ .env.regtest} (the
full stack, including the \texttt{boltz} profile) and the local nostr
relay. The first run pulls images and can take a few minutes; it's ready
when the CLI prints the service URLs (arkd on
\texttt{http://localhost:7070}, Fulmine on
\texttt{http://localhost:7003}, the Boltz CORS proxy on
\texttt{http://localhost:9069}, \ldots).

A handy alias for the counterparty Lightning node:

\begin{Shaded}
\begin{Highlighting}[]
\BuiltInTok{alias}\NormalTok{ lncli=}\StringTok{"docker exec {-}i lnd lncli {-}{-}network=regtest"}
\end{Highlighting}
\end{Shaded}

\hypertarget{start-the-wallet}{%
\subsection{2. Start the wallet}\label{start-the-wallet}}

\begin{Shaded}
\begin{Highlighting}[]
\ExtensionTok{pnpm}\NormalTok{ start}
\end{Highlighting}
\end{Shaded}

Open \url{http://localhost:3002} and create/unlock a wallet. On
\texttt{localhost} the app automatically targets the local \texttt{arkd}
(\texttt{http://localhost:7070}) and its regtest network --- no
\texttt{VITE\_ARK\_SERVER} needed. The regtest Boltz endpoint
(\texttt{http://localhost:9069}) is likewise built in.

\begin{quote}
To exercise the LNURL / nostr-backup features as well, start with:
\texttt{VITE\_LNURL\_SERVER\_URL=http://localhost:9090\ VITE\_NOSTR\_RELAY\_URL=ws://localhost:10547\ pnpm\ start}
\end{quote}

\hypertarget{fund-the-wallet}{%
\subsection{3. Fund the wallet}\label{fund-the-wallet}}

On the wallet's \textbf{Receive} screen, copy the on-chain (boarding)
address --- the \texttt{bcrt1…} one --- and faucet it:

\begin{Shaded}
\begin{Highlighting}[]
\CommentTok{\# 0.001 BTC; {-}{-}confirm mines a block so the deposit confirms}
\ExtensionTok{node}\NormalTok{ regtest/regtest.mjs faucet }\OperatorTok{\textless{}}\NormalTok{address}\OperatorTok{\textgreater{}}\NormalTok{ 0.001 }\AttributeTok{{-}{-}confirm}
\end{Highlighting}
\end{Shaded}

Back on the wallet home, open the pending transaction and
\textbf{settle} it. You're now ready to test swaps.

\hypertarget{test-submarine-swap-ark-lightning}{%
\subsection{Test Submarine Swap (Ark →
Lightning)}\label{test-submarine-swap-ark-lightning}}

Create a 5,000-sat invoice on the counterparty node and pay it from the
wallet:

\begin{Shaded}
\begin{Highlighting}[]
\ExtensionTok{lncli}\NormalTok{ addinvoice }\AttributeTok{{-}{-}amt}\NormalTok{ 5000 }\KeywordTok{|} \ExtensionTok{jq} \AttributeTok{{-}r}\NormalTok{ .payment\_request}
\end{Highlighting}
\end{Shaded}

Copy the \texttt{payment\_request} and pay it in Arkade. After payment
your transaction history shows the outgoing swap (amount + Boltz fee),
and Fulmine's history (\url{http://localhost:7003}) shows the matching
incoming movement.

\begin{quote}
The exact sat split (amount vs.~fee) depends on the Boltz release and
on-chain fee rate, so don't expect fixed numbers --- just that
\texttt{amount\ +\ fees} reconciles.
\end{quote}

\hypertarget{test-reverse-swap-lightning-ark}{%
\subsection{Test Reverse Swap (Lightning →
Ark)}\label{test-reverse-swap-lightning-ark}}

In Arkade go to \textbf{Receive}, enter an amount (e.g.~4,000 sats) and
continue to get a Lightning invoice. Pay it from the counterparty node:

\begin{Shaded}
\begin{Highlighting}[]
\ExtensionTok{lncli}\NormalTok{ payinvoice }\AttributeTok{{-}{-}force} \OperatorTok{\textless{}}\NormalTok{invoice}\OperatorTok{\textgreater{}}
\end{Highlighting}
\end{Shaded}

The wallet should show the received funds a moment later.

\hypertarget{teardown}{%
\subsection{Teardown}\label{teardown}}

\begin{Shaded}
\begin{Highlighting}[]
\ExtensionTok{pnpm}\NormalTok{ run regtest:stop     }\CommentTok{\# stop the stack}
\ExtensionTok{pnpm}\NormalTok{ run regtest:clean    }\CommentTok{\# stop and wipe all volumes/state}
\end{Highlighting}
\end{Shaded}

\hypertarget{troubleshooting}{%
\subsection{Troubleshooting}\label{troubleshooting}}

\begin{itemize}
\item
  On Apple Silicon (M-family) you may need to build the boltz-backend
  image locally and point the stack at it:

\begin{Shaded}
\begin{Highlighting}[]
\CommentTok{\# Clone boltz{-}backend locally}
\FunctionTok{git}\NormalTok{ clone git@github.com:BoltzExchange/boltz{-}backend.git }\KeywordTok{\&\&} \BuiltInTok{cd}\NormalTok{ boltz{-}backend}
\CommentTok{\# Build the image; VERSION=ark builds the ark branch.}
\ExtensionTok{docker}\NormalTok{ build }\AttributeTok{{-}{-}build{-}arg}\NormalTok{ NODE\_VERSION=lts{-}bookworm{-}slim }\AttributeTok{{-}{-}build{-}arg}\NormalTok{ VERSION=ark }\DataTypeTok{\textbackslash{}}
  \AttributeTok{{-}t}\NormalTok{ boltz/boltz:ark }\AttributeTok{{-}f}\NormalTok{ docker/boltz/Dockerfile .}
\end{Highlighting}
\end{Shaded}

  Then set \texttt{BOLTZ\_IMAGE=boltz/boltz:ark} in
  \texttt{.env.regtest} and re-run \texttt{pnpm\ run\ regtest:start}.
\end{itemize}

\end{document}
